\documentclass[11pt,a4paper]{article}

\usepackage[utf8]{inputenc}
\usepackage[T1]{fontenc}
\usepackage{lmodern}
\usepackage[margin=2.5cm]{geometry}
\usepackage{amsmath,amssymb,amsfonts}
\usepackage{booktabs}
\usepackage{xcolor}
\usepackage{enumitem}
\usepackage{titlesec}
\usepackage{fancyhdr}
\usepackage[most]{tcolorbox}
\usepackage[round,authoryear]{natbib}
\usepackage[hidelinks,colorlinks=true,linkcolor=blue!70!black,citecolor=blue!70!black,urlcolor=blue!70!black]{hyperref}

% Page styling
\setlength{\headheight}{14pt}
\pagestyle{fancy}
\fancyhf{}
\lhead{\textsc{Response to Reviewer 1 --- Dr.~Marcella Bonazzoli}}
\rhead{\textsc{Gerardo Cicalese}}
\cfoot{\thepage}

\definecolor{commentbg}{RGB}{245, 247, 250}
\definecolor{commentframe}{RGB}{70, 95, 130}
\definecolor{responseframe}{RGB}{40, 110, 70}

\newtcolorbox{revcomment}[2][]{
    enhanced,
    breakable,
    colback=commentbg,
    colframe=commentframe,
    fonttitle=\bfseries,
    title={#2},
    boxrule=0.8pt,
    arc=2mm,
    top=2mm, bottom=2mm, left=3mm, right=3mm,
    before skip=3mm, after skip=2mm,
    #1
}

\newenvironment{response}{
    \par\vspace{1mm}\noindent\textbf{\textcolor{responseframe}{Response:}}\par\vspace{1mm}
}{
    \par\vspace{3mm}
}

\begin{document}

\title{\textbf{\Large Formal Response to Reviewer 1 (Dr.~Marcella Bonazzoli)}\\[0.3em]
\large Re: Ph.D. Dissertation Review\\\emph{``Domain Decomposition and Differentiable Physics for the Viscoelastic Wave Equation''}}
\author{\textbf{Candidate:} Gerardo Cicalese\\[0.3em]
\textbf{Advisor:} Prof.~Ilario Mazzieri\\[0.3em]
\textbf{Co-Advisors:} Prof.~Gabriele Ciaramella, Prof.~Fabio Antonacci\\[0.3em]
\textbf{Institution:} Politecnico di Milano}
\date{November 2026}
\maketitle

\noindent Dear Dr.~Marcella Bonazzoli,\\[0.5em]
We would like to express our sincere gratitude for your thorough, encouraging, and constructive evaluation of the doctoral dissertation \emph{``Domain Decomposition and Differentiable Physics for the Viscoelastic Wave Equation''}. We are deeply honored by your positive assessment of the originality, theoretical rigor, and balanced interplay between analytical derivations and extensive numerical simulations presented in the manuscript.

Your detailed comments and precise mathematical observations have been invaluable in refining the manuscript. They have allowed us to eliminate residual ambiguities, correct formal mathematical definitions, refine transmission operator descriptions, provide complete algorithmic details, and ensure absolute consistency between our analytical statements and numerical figures.

In the revised dissertation, modifications specifically addressing your comments are highlighted in \textcolor{blue}{blue}; modifications addressing comments shared across the reviewer reports are highlighted in \textcolor{red}{red}. This distinction is used throughout the manuscript (\texttt{\textbackslash textcolor\{blue\}\{...\}} and \texttt{\textbackslash textcolor\{red\}\{...\}}).

Below, we provide a point-by-point response to your General Assessment and all 19 detailed comments and suggestions.

\tableofcontents
\newpage

% ==============================================================================
\section{Response to the General Assessment}
% ==============================================================================

\begin{revcomment}{General Assessment by Reviewer 1}
\emph{``I have found the contents of the PhD thesis `Domain Decomposition and Differentiable Physics for the Viscoelastic Wave Equation' by Gerardo Cicalese very interesting, and reading the manuscript has been a rewarding experience: the various topics covered by the thesis are clearly explained, the needed background is correctly recalled, and the novelties contained in each chapter are easy to detect. The thesis is a balanced interplay between theory and simulations: all the theoretical results are validated by extensive numerical experiments. [...] It is remarkable that the candidate was able to successfully and significantly contribute to the knowledge advance in three research domains that are often studied separately: stable discretization methods, efficient space-time solvers, and neural networks for inverse problems. As a general assessment, I consider the thesis as a very significant and excellent contribution to the knowledge advance in a field where several outstanding researchers are working actively. In view of the above considerations, I believe that the PhD thesis by Gerardo Cicalese conforms to the standards of a high quality international thesis, and it is my pleasure to recommend that he be admitted to the final defense of his doctoral dissertation.''}
\end{revcomment}

\begin{response}
We are profoundly grateful to Dr.~Marcella Bonazzoli for this exceptionally generous, encouraging, and positive appraisal of the doctoral dissertation. Her recognition of the novel contributions across all three pillars of the research---stable numerical discretizations for dissipative media, space-time domain decomposition methods, and differentiable physics architectures for inverse acoustics---is a tremendous encouragement.

In revising the manuscript, we have addressed every one of her 19 detailed comments with the utmost care, ensuring mathematical precision and clarity throughout the document.
\end{response}

\newpage
% ==============================================================================
\section{Point-by-Point Responses to Detailed Comments}
% ==============================================================================

% ------------------------------------------------------------------------------
\subsection{Comments on Chapter 1: Introduction and Literature Review}
% ------------------------------------------------------------------------------

\begin{revcomment}{Comment 1 (Page 9): Typo in Room Acoustic Methods}
\emph{``Page 9: what are the `tan' methods?''}
\end{revcomment}

\begin{response}
We thank the reviewer for spotting this clerical typo. The stray ``tan'' text was removed. The current introduction states: \textcolor{blue}{\emph{``Although time-dependent wave-based methods date back to the late 1960s and 1970s in geophysics and seismology, their application to room acoustics emerged in the 1990s.''}} This wording distinguishes the development of wave-based schemes from their later adoption in room acoustics.
\end{response}

% ------------------------------------------------------------------------------
\subsection{Comments on Chapter 2: Physical Model and Discretization Methods}
% ------------------------------------------------------------------------------

\begin{revcomment}{Comment 2 (Pages 18, 21): Precision on Robin / DtN Transparent Boundary Conditions}
\emph{``Pages 18, 21: the sentences `in the absence of damping, the Robin boundary condition coincides with the exact transparent boundary condition', `the reflection vanishes if $\Lambda(i\omega) = ik(i\omega)$, that is, when the Robin operator coincides with the exact Dirichlet-to-Neumann map' are not precise. Indeed, the impedance choice $i\omega/c$ is a zero-th order approximation of the exact DtN transmission operator, see e.g. [\S2.5.2, Dolean, Jolivet, Nataf 2015].''}
\end{revcomment}

\begin{response}
We completely agree with the reviewer's observation. The exact continuous Dirichlet-to-Neumann (DtN) operator for the wave equation in multi-dimensional domains or for oblique incident waves is non-local in space and time, depending intrinsically on the transverse wavenumber $k_\perp$ via the exact symbol $\sqrt{\omega^2/c^2 - |k_\perp|^2}$. The local Robin impedance choice $\Lambda(i\omega) = i\omega/c$ ($q=1/c, r=0$) coincides with the exact transparent DtN condition \emph{strictly and exclusively for one-dimensional plane waves propagating at normal incidence}. In general multi-dimensional settings, it constitutes a local zeroth-order Taylor/Engquist--Majda approximation to the exact DtN map.

We have revised both passages in the manuscript accordingly, citing the recommended monograph by \citet{dolean_introduction_2015} as well as the foundational work of \citet{engquist_absorbing_1977}:
\begin{itemize}[leftmargin=*]
    \item \textbf{Section~2.1.2:}
    \begin{quote}
    \textcolor{blue}{\emph{``Therefore, in the absence of damping, the Robin boundary condition coincides with the exact Dirichlet-to-Neumann (DtN) transparent boundary condition for one-dimensional plane waves at normal incidence, yielding perfect absorption across all frequencies. In multiple spatial dimensions, the exact DtN operator is non-local and angle-dependent; hence, the local Robin condition (2.2) represents a zeroth-order local approximation (in the sense of Engquist and Majda \cite{engquist_absorbing_1977,dolean_introduction_2015}) to the true transparent boundary condition.''}}
    \end{quote}
    \item \textbf{Section~2.3.1:}
    \begin{quote}
    \textcolor{blue}{\emph{``that is, when the Robin operator coincides with the exact Dirichlet-to-Neumann map (which provides exact transparent absorption for 1D plane waves at normal incidence, but represents a zeroth-order local approximation in higher-dimensional geometries \cite{dolean_introduction_2015}).''}}
    \end{quote}
\end{itemize}
\end{response}

\vspace{2mm}
\begin{revcomment}{Comment 3 (Page 27): Forward Definition of Eigenfunctions $\phi_m(x)$}
\emph{``Page 27: the $\phi_m$ have not been defined yet.''}
\end{revcomment}

\begin{response}
We thank the reviewer for highlighting this forward reference. In Section~2.4.1 (Configuration~I), the spatial modal expansion $u(x,t) = \sum_{m \in \{1,3\}} U_m(t)\phi_m(x)$ was introduced prior to the general modal analysis in Section~2.7.

In the revised manuscript, we have added the explicit definition directly at its first occurrence in Configuration~I:
\begin{quote}
\textcolor{blue}{\emph{``where $\phi_m(x) = \cos(m\pi x / L)$ denotes the spatial eigenfunctions of the Laplacian subject to homogeneous Neumann boundary conditions.''}}
\end{quote}
\end{response}

\vspace{2mm}
\begin{revcomment}{Comments 4 \& 5 (Pages 30--31): Scheme (2.8) and Redundant Non-Centered Formulation}
\emph{``Page 30: $f_i^n$ is missing in (2.8).''}

\emph{``Pages 30-31: scheme (2.8) is actually never used. Since the derivation of scheme (2.9) (whose only difference with (2.8) is the centered approximation for $\partial_t$) is much more detailed, I suggest removing scheme (2.8).''}
\end{revcomment}

\begin{response}
We fully agree with the reviewer's suggestion. Scheme (2.8) in the previous draft was an introductory sketch that employed a backward-difference approximation in time for the mixed viscoelastic term, omitted the external force $f_i^n$, and was never actually implemented or analyzed in the subsequent stability proofs or numerical experiments.

Following the reviewer's recommendation, \textbf{we have completely removed the former introductory scheme (2.8)} from Section~2.5. The section now proceeds directly from the space and time grid discretization into the derivation of the centered leapfrog FDTD scheme, now numbered Equation~(2.8), which properly includes the discrete source term $f_i^n$ and centered second-order approximations for the temporal and spatial derivatives.
\end{response}

\vspace{2mm}
\begin{revcomment}{Comment 6 (Page 49): Explicit Expression of the DCT-II Matrix $\Phi$}
\emph{``Page 49: it would be useful to give the explicit expression of matrix $\Phi$.''}
\end{revcomment}

\begin{response}
We agree that writing out the explicit matrix formula is helpful for the reader. In Section~2.7.2, we have explicitly specified the transformation matrix $\Phi \in \mathbb{R}^{N \times N}$ in terms of the discrete orthonormal DCT-II basis functions:
\begin{quote}
\textcolor{red}{
\emph{``For convenience, we introduce the orthonormal DCT-II change-of-basis matrix $\Phi \in \mathbb{R}^{N \times N}$, whose $(i,m)$-th entry is given by the discrete basis functions:
\begin{equation*}
\Phi_{i,m} = \phi_m(i) = \sqrt{\frac{2}{N}} \alpha_m \cos\!\left(\frac{m\pi}{N}\left(i + \frac{1}{2}\right)\right), \quad i, m = 0, \dots, N-1,
\end{equation*}
with scaling factors $\alpha_0 = 1/\sqrt{2}$ and $\alpha_m = 1$ for $m \ge 1$, satisfying $\Phi^{-1} = \Phi^T$.''}
}
\end{quote}
\end{response}

\vspace{2mm}
\begin{revcomment}{Comment 7 (Page 50): Definition of Permuted Vectors and Permutation Matrix $P$}
\emph{``Page 50: the definition of the permuted versions $\widetilde{\mathbf{Z}}, \widetilde{\mathbf{W}}$ is not clear, which permutation is precisely matrix $P$?''}
\end{revcomment}

\begin{response}
We thank the reviewer for requesting clarification on this algebraic step. In the modal state-space formulation, the $2N$-dimensional state vector naturally arises in pairwise interleaved form, grouping modal displacement and velocity for each individual mode $m$:
\[
\mathbf{Z} = [U_0, V_0, U_1, V_1, \dots, U_{N-1}, V_{N-1}]^T \in \mathbb{R}^{2N}.
\]
However, applying the orthogonal block Discrete Cosine Transform $\operatorname{diag}(\Phi, \Phi)$ requires grouping all modal displacements together followed by all modal velocities:
\[
\widetilde{\mathbf{Z}} = [\mathbf{U}^T, \mathbf{V}^T]^T = [U_0, U_1, \dots, U_{N-1}, V_0, V_1, \dots, V_{N-1}]^T \in \mathbb{R}^{2N}.
\]
In the revised manuscript (Section~2.7.2), we have added an explicit explanation:
\begin{quote}
\textcolor{red}{\emph{``where $P$ is the $2N\times2N$ orthogonal permutation matrix ($P^{-1} = P^T$) that reorders the pairwise interleaved modal state vector
\begin{equation*}
\mathbf Z = [U_0, V_0, U_1, V_1, \dots, U_{N-1}, V_{N-1}]^T
\end{equation*}
into the segregated block-component vector $\widetilde{\mathbf Z} = [\mathbf U^T, \mathbf V^T]^T$ grouping all modal displacements first, followed by all modal velocities.''}}
\end{quote}
\end{response}

\vspace{2mm}
\begin{revcomment}{Comment 8 (Page 55): Bibliographic References for Fast DCT Algorithms}
\emph{``Page 55: can you give bibliographic references for fast DCT algorithms?''}
\end{revcomment}

\begin{response}
We have added foundational and standard references for fast $O(N \log N)$ Discrete Cosine Transform algorithms in Section~2.7.4:
\begin{itemize}[leftmargin=*]
    \item \citet{makhoul_fast_1980}: \emph{``A fast cosine transform in one and two dimensions''}, IEEE Transactions on Acoustics, Speech, and Signal Processing;
    \item \citet{rao_discrete_1990}: \emph{``Discrete Cosine Transform: Algorithms, Advantages, Applications''}, Academic Press;
    \item \citet{frigo_design_2005}: \emph{``The Design and Implementation of FFTW3''}, Proceedings of the IEEE.
\end{itemize}
The text now reads:
\begin{quote}
\textcolor{red}{\emph{``forward and inverse transforms are evaluated using fast DCT algorithms with computational complexity $O(N\log N)$ \cite{makhoul_fast_1980,rao_discrete_1990,frigo_design_2005}.''}}
\end{quote}
\end{response}

% ------------------------------------------------------------------------------
\subsection{Comments on Chapter 3: Domain Decomposition and Inverse Problems}
% ------------------------------------------------------------------------------

\begin{revcomment}{Comment 9 (Page 75): Solution of the OSWR Min-Max Optimization Problem}
\emph{``Page 75: how is the min-max problem (3.1) solved? Do you have an explicit expression of the optimal parameters $q, r$?''}
\end{revcomment}

\begin{response}
We thank the reviewer for raising this methodological point. We state the continuous-frequency min--max objective, but select the parameters numerically over a sampled frequency set:
\[
\min_{q, r \ge 0} \max_{\omega \in [\omega_{\min}, \omega_{\max}]} \rho(\omega; q, r)
\]
For arbitrary $(\gamma,\nu)$, this continuous-frequency problem has no general closed-form analytical solution.

In Section~3.1.2 and the expanded Future Work section (Section~4.3.3), we have explicitly clarified:
\begin{enumerate}[leftmargin=*]
    \item The parameters are selected using MATLAB's constrained \texttt{patternsearch} optimizer over 4,000 logarithmically spaced frequency samples, with $q,r\ge0$ and initial guess $(1/c,\gamma/(2c))$. The minimized quantity is $|G|^2$, the contraction factor over the two-iteration recurrence $A_u^{(k+2)}=G^2A_u^{(k)}$. This is a numerical sampled-grid solution, so it does not rule out a larger value between samples.
    \item Closed-form expressions for the combined telegrapher--viscoelastic case are left for future work, as now stated in Section~4.3.3.
\end{enumerate}
\end{response}

\vspace{2mm}
\begin{revcomment}{Comment 10 (Page 78): Citation of Figure 3.3 in Main Text}
\emph{``Page 78: Figure 3.3 is not mentioned in the main text.''}
\end{revcomment}

\begin{response}
We thank the reviewer for catching this omission. In Section~3.1.3, we have added an explicit citation and descriptive discussion of Figure~3.3:
\begin{quote}
    \textcolor{blue}{\emph{``Figure~3.3 displays the parameters $q^\ast$ and $r^\ast$ selected by the sampled-frequency min--max search, together with the corresponding predicted factor $\rho_{\mathrm{opt}}$. In the tested parameter range, the sampled-frequency factors remain below one. Since the search uses a finite frequency set, this result does not establish a bound on the continuous-frequency maximum.''}}
\end{quote}
\end{response}

\vspace{2mm}
\begin{revcomment}{Comment 11 (Page 80): Removal of Unannotated Scatter Plot Figure 3.5}
\emph{``Page 80: Figure 3.5 is not very informative because it does not give values for $\gamma, \nu$. Since Figure 3.6 does give values for $\gamma, \nu$, maybe you could remove Figure 3.5.''}
\end{revcomment}

\begin{response}
We fully agree with the reviewer's recommendation. Figure~3.5 in the previous draft was an unannotated scatter plot aggregating all $(q, r)$ points without identifying the underlying physical damping parameters $\gamma$ and $\nu$. It has been removed. The revised Figure~3.5 shows parametric trajectories along constant $\gamma$ and constant $\nu$, with the damping values indicated by color.

Following the reviewer's suggestion, we removed the unannotated scatter plot and updated the surrounding text to introduce the parametric trajectory curves directly.
\end{response}

\vspace{2mm}
\begin{revcomment}{Comment 12 (Page 81): Explicit Definition of Domain $\Omega$ in 2D OSWR}
\emph{``Page 81: please define explicitly the domain $\Omega$ for the 2D experiment.''}
\end{revcomment}

\begin{response}
We thank the reviewer for noting this omission. In Section~3.1.4, we have updated the opening sentence to explicitly define the global two-dimensional computational domain $\Omega$:
\begin{quote}
\textcolor{red}{\emph{``Let us consider the damped wave equation in two spatial dimensions $(x, y)$ on the rectangular domain $\Omega = (0, L_x) \times (0, L_y)$, where the domain is partitioned along the $x$-axis (allowing the transverse $y$-direction to be decoupled via a spatial Fourier transform) into two overlapping subdomains $\Omega_1 = (0, b) \times (0, L_y)$ and $\Omega_2 = (a, L_x) \times (0, L_y)$ with $0 < a < b < L_x$ (sharing the overlap width $\delta = b - a > 0$).''}}
\end{quote}
\end{response}

\vspace{2mm}
\begin{revcomment}{Comment 13 (Page 92): Convergence Rate in ARD Interface Operator (Figure 3.9)}
\emph{``Page 92: it is written that `the domain decomposition error converges linearly', but in the legend of the corresponding Figure 3.9 you have put quadratic convergence.''}
\end{revcomment}

\begin{response}
We are very grateful to the reviewer for pointing out this apparent contradiction between the text and the figure legend. In the previous draft, Figure~3.9 plotted only an $O(\Delta x^2)$ reference line, while the accompanying text stated that the convergence was linear ($O(\Delta x)$).

To resolve this issue rigorously:
\begin{enumerate}[leftmargin=*]
    \item \textbf{Updated Convergence Figures:} We have re-executed the grid refinement tests and updated all six subfigures of Figure~3.8 to display \textbf{both} a quadratic reference line ($O(\Delta x^2)$, dashed black line) and a linear reference line ($O(\Delta x)$, dotted black line) with clear legends.
    \item \textbf{Physical and Numerical Interpretation:} The revised text in Section~3.2.4 analyzes the observed behavior against both reference slopes:
    \begin{quote}
    \textcolor{red}{\emph{``As observed in Figure~3.8, which displays quadratic ($O(\Delta x^2)$) and linear ($O(\Delta x)$) reference slopes, the empirical domain decomposition error exhibits different trends across the tested profiles and grid refinements. For the smooth standing-wave configurations (a) and (b), the error decreases approximately quadratically on coarser grids and approaches a linear trend on finer meshes. For the non-smooth configurations (c, d, f) and propagating pulses, the error is predominantly consistent with a linear trend over the tested resolutions. These observations are compatible with contributions from interface truncation and splitting error, but the experiments do not isolate their individual effects or establish a general asymptotic rate.''}}
    \end{quote}
\end{enumerate}
\end{response}

\vspace{2mm}
\begin{revcomment}{Comment 14 (Page 96): 3D Domain Boundary Conditions, Damping Parameter Count, and Viscosity Differences}
\emph{``Page 96: Please specify whether the original open environment is unbounded in the $z$ direction or also in other directions. When you write `three different viscoelastic damping coefficients', it is actually two. You say `the progressive attenuation and smoothing of the wavefronts scale consistently with the imposed viscoelastic parameter $\nu$', but in Figure 3.12 the difference between plots (a)-(b)-(c) with $\nu = 2 \cdot 10^{-6}$ and plots (d)-(e)-(f) with $\nu = 10^{-5}$ is not clear visually.''}
\end{revcomment}

\begin{response}
We thank the reviewer for these precise observations. We have addressed the boundary description and damping comparison in the Chapter~3 3D ARD experiment discussion:
\begin{enumerate}[leftmargin=*]
    \item \textbf{Boundary Conditions:} We explicitly clarified that the 3D room simulation is enclosed by rigid bounding walls with homogeneous Neumann boundary conditions on all six outer faces ($x, y, z$), which allows the exact DCT modal representation to be employed within each partitioned rectangular subdomain.
    \item \textbf{Damping settings:} The revised snapshots compare four values: an undamped case ($\nu=0$) and three positive values ($\nu\approx2.4\times10^{-6}$, $1.2\times10^{-5}$, and $5.9\times10^{-5}$\,m$^2$/s).
    \item \textbf{Scope of the comparison:} The snapshots provide a qualitative comparison of the displayed fields. The receiver frequency-response comparison is for the lossless case and focuses on cross-solver consistency below $200$\,Hz; it is not a quantitative measurement of attenuation as a function of $\nu$. The revised text therefore avoids claiming a spectral verification of the damping sweep.
\end{enumerate}
\end{response}

\vspace{2mm}
\begin{revcomment}{Comment 15 (Pages 97, 101): Sensor Measurement Operator Details}
\emph{``Page 97 or page 101: please give more details about the considered sensor measurements over time (are their supports the points $x_i$? Do you have measurements at each time step?)''}
\end{revcomment}

\begin{response}
We thank the reviewer for requesting complete technical details on the measurement setup. In Sections~3.3.1 and~3.3.2, we have provided the mathematical definition of the discrete observation operator $\mathcal{M}$:
\begin{itemize}[leftmargin=*]
    \item The benchmark uses 10 fixed interior PSTD cell-centered receivers at $x_i=(i+\tfrac12)\Delta x$, selected once using seed 123. For FDTD, each location is mapped to the nearest endpoint-grid node.
    \item Measurements are recorded at all 41 time samples $t_n=n\Delta t$, $n=0,\ldots,40$, for $T=2$\,s and $\Delta t=0.05$\,s; there is no temporal subsampling.
    \item The observation array is in $\mathbb{R}^{10\times41}$. Additive Gaussian noise is normalized to 30\,dB relative SNR; the seed is reset for each profile.
\end{itemize}
\end{response}

\vspace{2mm}
\begin{revcomment}{Comment 16 (Page 98): Parameter Mapping and Softplus References}
\emph{``Page 98: can you give bibliographic references for Softplus mappings?''}
\end{revcomment}

\begin{response}
In the revised methodology (Section~3.3.1), we replaced the positive-activation mapping with a logarithmic parameterization:
\[
\gamma = \gamma_{\mathrm{ref}} e^{\xi_\gamma}, \qquad \nu = \nu_{\mathrm{ref}} e^{\xi_\nu},
\]
This mapping enforces positive parameter values. The revised formulation does not use a Softplus mapping, so the Softplus citation is not part of the parameterization used in the thesis.
\end{response}

\vspace{2mm}
\begin{revcomment}{Comment 17 (Page 99): Regularization Scale and Cost Functional}
\emph{``Page 99: is the regularization scale $\nu$ in the cost functional the $\widetilde{\nu}$ in its argument?''}
\end{revcomment}

\begin{response}
In the original draft, the proposed loss function included an $H^1$ Sobolev smoothness penalty scaled by the viscoelastic parameter $\nu$. As noted by both reviewers, coupling the loss functional directly to an unknown optimization variable introduces complex scaling ambiguities and unintended feedback during gradient updates.

In the revised dissertation (Section~3.3.1), \textbf{we have eliminated the $H^1$ penalty from the proposed DPNN loss function}. The primary DPNN formulation now operates as a pure data-fitting objective:
\[
\mathcal{L}_{\mathrm{data}}(\Theta, \xi_\gamma, \xi_\nu) = \frac{1}{2 S N_t} \sum_{s=1}^S \sum_{n=0}^{N_t} \left| u(x_s, t_n; \Theta, \xi_\gamma, \xi_\nu) - u_{\mathrm{obs}}(x_s, t_n) \right|^2.
\]
The $H^1$ and $L^1$ penalties are retained within the ablation study (Section~3.3.4, Table~3.5). Among the tested configurations, data-only fitting gives the better combined field-and-parameter trade-off; $L^1$ lowers the Gaussian initial-condition NMSE but worsens its field error, and the tested $H^1$ configurations perform poorly.
\end{response}

\vspace{2mm}
\begin{revcomment}{Comment 18 (Page 102): Physical Parameter Convergence ($\gamma$ vs $\nu$)}
\emph{``Page 102: you write that the physical parameters $\gamma$ and $\nu$ converge to `values close to the true ones', but in Figure 3.14 (and Figures 3.17, 3.20) it appears that $\gamma$ is not well reconstructed actually.''}
\end{revcomment}

\begin{response}
We thank the reviewer for this critical observation. The revised results show that parameter recovery depends on the initial profile: PSTD-DPNN has a large $\gamma$ error in the multimodal standing-wave case, while the reported Gaussian and triangular runs recover both parameters more accurately. The estimates therefore do not support a blanket statement that $\gamma$ is poorly identified or that $\nu$ is consistently recovered with high precision across all benchmarks.

In the revised dissertation, we have completely overhauled the discussion in Section~3.3.2, Section~3.3.3, and the Conclusion (Section~4.1) to state this result honestly and rigorously:
\begin{enumerate}[leftmargin=*]
    \item We report the profile-dependent parameter errors and limit the conclusion to the tested profiles, sensors, noise level, and training runs;
    \item We provide a Gaussian parameter-only run and a Gaussian joint run in Section~3.3.3 (Table~3.4), but these results do not establish a general source-damping coupling mechanism:
    \begin{itemize}
        \item In the parameter-only Gaussian run, the reported relative errors are $10.8\%$ for $\gamma$ and $3.0\%$ for $\nu$; in the joint Gaussian run they are $3.5\%$ and $2.0\%$, respectively. These single-run values do not show joint estimation to be consistently harder.
        \item The results motivate further investigation of source--parameter confounding, but do not establish that it caused the profile-dependent errors observed here.
        \item The viscoelastic term has wavenumber-dependent modal effects in the model, but the experiments do not establish that this is the cause of the observed parameter-recovery differences.
    \end{itemize}
\end{enumerate}
\end{response}

\vspace{2mm}
\begin{revcomment}{Comment 19 (Page 104): Visual Comparison of PSTD vs FDTD Reconstruction}
\emph{``Page 104: in the caption of Figure 3.19 you write `Notice the superior amplitude preservation in the PSTD reconstruction', but visually the result with FDTD looks nearly the same as the PSTD one in this figure.''}
\end{revcomment}

\begin{response}
We agree with the reviewer that claiming an obvious visual difference in the initial-profile comparison was an overstatement. The reviewer refers to Figure~3.19 in the previous draft; the corresponding Gaussian initial-condition comparison is Figure~3.20 in the revised thesis.

In the revised caption and main text (Section~3.3.2), we have toned down this claim:
\begin{quote}
    \textcolor{red}{\emph{``Figure~3.20 shows that both DPNN solvers capture the Gaussian envelope; Table~3.2 provides the quantitative comparison.''}}
\end{quote}
We emphasize the quantitative comparisons in Table~3.2. The modal propagator has no temporal integration error for retained modes, while spatial truncation and representation errors remain; the reported experiments do not establish zero dispersion for all modes or a general long-horizon advantage.
\end{response}

\newpage
% ==============================================================================
\section{Summary of Manuscript Changes}
% ==============================================================================

\begin{table}[ht]
\centering
\small
\caption{Comprehensive Summary of Revisions in Response to Reviewer 1 (Dr.~Marcella Bonazzoli).}
\vspace{2mm}
\begin{tabular}{@{}llp{8.5cm}l@{}}
\toprule
\textbf{Comment} & \textbf{Section} & \textbf{Nature of Revision} & \textbf{Status} \\
\midrule
1 & \S1.2, \S1.4 & Corrected typo: ``tan'' methods replaced by wave-based time-domain methods & Addressed \\
2 & \S2.1.2, \S2.3.1 & Clarified that Robin impedance $i\omega/c$ is exact DtN for 1D normal incidence and a zeroth-order local approximation in multi-D \citep{dolean_introduction_2015} & Addressed \\
3 & \S2.4.1 & Added explicit forward definition of spatial eigenfunctions $\phi_m(x) = \cos(m\pi x/L)$ in Configuration~I & Addressed \\
4, 5 & \S2.5 & Removed the former introductory scheme (2.8); derived the centered leapfrog scheme with $f_i^n$ & Addressed \\
6 & \S2.7.2 & Provided explicit matrix formula for orthonormal DCT-II matrix $\Phi$ & Addressed \\
7 & \S2.7.2 & Clarified permutation matrix $P$ reordering interleaved $[U_m, V_m]$ into segregated block vector & Addressed \\
8 & \S2.7.4 & Added standard citations for fast DCT algorithms \citep{makhoul_fast_1980,rao_discrete_1990,frigo_design_2005} & Addressed \\
9 & \S3.1.2, \S4.3.3 & Specified the constrained pattern-search minimization over sampled frequencies and the absence of a general closed-form solution & Addressed \\
10 & \S3.1.3 & Added in-text citation and discussion of Figure~3.3 & Addressed \\
11 & \S3.1.3 & Removed the former unannotated scatter plot; retained parametric trajectories as Figure~3.5 & Addressed \\
12 & \S3.1.4 & Explicitly defined the 2D global domain $\Omega = (0, L_x) \times (0, L_y)$ & Addressed \\
13 & \S3.2.4 & Regenerated Figure~3.8 with both $O(\Delta x)$ and $O(\Delta x^2)$ reference lines; qualified the observed trends & Addressed \\
14 & 3D ARD experiment & Clarified the rigid outer boundary conditions and four damping values; distinguished qualitative damping snapshots from the lossless cross-solver frequency-response comparison & Addressed \\
15 & \S3.3.1, \S3.3.2 & Specified PSTD cell-center coordinates, nearest-node FDTD mapping, and all 41 time samples & Addressed \\
16 & \S3.3.1 & Replaced Softplus with a logarithmic positive-parameter mapping & Addressed \\
17 & \S3.3.1, \S3.3.4 & Removed $H^1$ regularization from proposed DPNN loss; retained it in the comparative ablation study & Addressed \\
18 & \S3.3.2, \S3.3.3 & Reported profile-dependent parameter errors and limited the source--damping interpretation to a hypothesis for future study & Addressed \\
19 & \S3.3.2 & Toned down visual superiority claims in Figure~3.20; referred readers to the quantitative results in Table~3.2 & Addressed \\
\bottomrule
\end{tabular}
\end{table}

\noindent In conclusion, we are deeply grateful to Dr.~Marcella Bonazzoli for her thorough and constructive review, which has significantly enhanced the precision, clarity, and scholarly quality of this doctoral dissertation.

\vspace{1cm}
\noindent Respectfully submitted,\\[1em]
\textbf{Gerardo Cicalese}\\
Ph.D. Candidate in Information Technology, Politecnico di Milano\\[0.5em]
\textbf{Prof.~Ilario Mazzieri}\\
Advisor, Associate Professor, Politecnico di Milano\\[0.5em]
\textbf{Prof.~Gabriele Ciaramella}\\
Co-Advisor, Associate Professor, Politecnico di Milano\\[0.5em]
\textbf{Prof.~Fabio Antonacci}\\
Co-Advisor, Associate Professor, Politecnico di Milano

\newpage
% ==============================================================================
% References
% ==============================================================================
\bibliographystyle{unsrtnat-nourl}
\bibliography{Thesis_bibliography}

\end{document}
